% titel: Blatt 5, A2: Grenzwerte von Funktionen
% Teil II: der durchgerechnete Weg.

\subsection*{Gegeben und gesucht}
Fünf Grenzwerte, die die drei Typen (Polynomverhältnis am Unendlichen,
einseitiger Grenzwert, $\infty-\infty$-Differenz) abdecken.

\subsection*{a) $\displaystyle\lim_{x\to\infty}\frac{-2x-3}{1+2x}$}
\emph{Gedankenschritt:} Zähler und Nenner sind beide Grad $1$. Wir kürzen
durch die höchste Potenz $x$; die konstanten Restterme $\tfrac3x$ bzw.
$\tfrac1x$ streben gegen $0$ und bleiben nur als Störung stehen.
\begin{equation*}
  \frac{-2x-3}{1+2x}=\frac{-2-\frac3x}{\frac1x+2}
\;\Rightarrow\;\frac{-2}{2}=-1.
\end{equation*}
Plausibel: Für sehr große $x$ ist der Bruch $-2x/(2x)=-1$; die kleinen
Korrekturen verschwinden.

\subsection*{b) $\displaystyle\lim_{x\to\infty}\sqrt{\frac{2x^2-x}{3x+8x^2}}$}
\emph{Gedankenschritt:} Innerhalb der Wurzel liegt ein Quotient vom Typ
\emph{gleicher Grad} ($2x^2$ zu $8x^2$). Erst den Bruch vereinfachen, dann die
(stetige) Wurzel ziehen.
\begin{equation*}
  \frac{2x^2-x}{3x+8x^2}=\frac{2-\frac1x}{\frac3x+8}
  \;\Rightarrow\;\frac{2}{8}=\frac14 .
\end{equation*}
Wegen der Stetigkeit von $\sqrt{\cdot}$:
$\lim\sqrt{(\cdot)}\ =\ \sqrt{\lim(\cdot)}\ =\ \sqrt{\frac14}=\frac12$.

\subsection*{c) $\displaystyle\lim_{x\to\infty}\ln\!\frac{2x^2-1}{2x+2x^2}$}
\emph{Gedankenschritt:} Man darf die Argumente von $\ln$ erst vereinfachen.
Kürzen durch die höchste Potenz $2x^2$ (Zähler und Nenner vom Grad $2$):
\begin{equation*}
  \frac{2x^2-1}{2x+2x^2}=\frac{2-\frac1{x^2}}{\frac1x+2}
  \;\Rightarrow\;\frac{2}{2}=1 .
\end{equation*}
Nun $\ln$ auf den Grenzwert anwenden (Stetigkeit): $\ln 1=0$.

\subsection*{d) $\displaystyle\lim_{x\searrow 1}\frac{2-(x+1)^3}{(x-1)^3}$}
\emph{Gedankenschritt:} Einsetzen $x=1$: Zähler $=2-8=-6$, Nenner $=0$.
Das ist ein Quotient \emph{endlich durch Null}, also divergiert er; das
Vorzeichen entscheidet die Richtung. Für $x\searrow 1$ (von rechts) ist
$x-1>0$, also $(x-1)^3>0$, der Zähler bleibt $-6<0$: Der Bruch geht gegen
$-\infty$.
\begin{equation*}
  \lim_{x\searrow 1}\frac{2-(x+1)^3}{(x-1)^3}=-\infty .
\end{equation*}

\subsection*{e) $\displaystyle\lim_{x\to\infty}\left(x-\sqrt{x^2-1}\right)$}
\emph{Gedankenschritt:} Es ist von der Form $\infty-\infty$, was
unbestimmt ist. Die Wurzel lässt sich durch Erweitern mit der konjugierten
Wurzel eliminieren:
\begin{equation*}
  x-\sqrt{x^2-1}
  =\frac{(x-\sqrt{x^2-1})(x+\sqrt{x^2-1})}{x+\sqrt{x^2-1}}
  =\frac{x^2-(x^2-1)}{x+\sqrt{x^2-1}}
  =\frac{1}{x+\sqrt{x^2-1}}\;\Rightarrow\;0 .
\end{equation*}
Plausibel: Für große $x$ ist $\sqrt{x^2-1}$ fast so groß wie $x$, und die
Differenz schrumpft auf den Wert $0$ (der Restterm ist $\frac{1}{2x}$).

\begin{ergebnis}
\begin{equation*}
  \text{a) } -1\;,\quad \text{b) } \tfrac12\;,\quad \text{c) } 0\;,\quad
  \text{d) } -\infty\;,\quad \text{e) } 0 .
\end{equation*}
\end{ergebnis}

\begin{kontrolle}
\textbf{Numerische Probe:} Für $x=10^6$ liefert a) $-1{,}000003$, b)
$0{,}4999992$, c) $-0{,}000001$, e) $1/(10^6+10^6)\approx5\cdot10^{-7}$ --
alles plausibel. Für d) mit $x=1{,}001$: Nenner $=10^{-9}$, Zähler $-6$,
Quotient $\approx -6\cdot10^9$, also \emph{riesig negativ}. \checkmark

\textbf{Grenzfall d:} Für $x\swarrow 1$ (von links) ist
$(x-1)^3<0$, also ginge der Bruch gegen $+\infty$; die Richtung hängt vom
Vorzeichen des Nenners ab. Dieses Umschlagen der Richtung ist typisch für
einseitige Grenzwerte. \checkmark
\end{kontrolle}

\begin{merksatz}
Am Unendlichen dominiert der höchste Grad; Zähler und Nenner durch die
höchste Potenz kürzen. Für $\infty-\infty$ mit Wurzeln mit der konjugierten
Wurzel erweitern. Stetige Funktionen wie $\sqrt{\cdot}$ und $\ln$
vertauschen mit dem Grenzwert, sobald das Argument konvergiert. Einseitige
Grenzwerte hängen getrennt vom Vorzeichen der Umgebung ab.
\end{merksatz}
